\section{A canonical positive count for order-ideal jet extensions}

\subsection{O-024: subtracting the jet order from every suffix bound}
Let $\lambda$ have length at most $r\geq2$, and let
$b\geq0$ be an integer. Set
\[
           \mu_j=(\lambda_j-b)_+\quad(2\leq j\leq r),
            \qquad u_+=\max(u,0),
\]
and write $\Gamma_\lambda$ for the full nested suffix
support of M-018. Then
\[
                         \Gamma_\lambda\subseteq\Gamma_\mu+D_b.
                                                               \tag{O-024.1}
\]
Moreover
\[
        |\Gamma_\mu|
            =\#\{\nu:\nu\text{ is a partition and }
                         0\leq\nu_j\leq\mu_j\ (2\leq j\leq r)\}.
                                                               \tag{O-024.2}
\]
Thus the needed extension size has a direct positive count:
the number of subpartitions of the tail obtained by
subtracting $b$ and replacing negative entries by zero.
This is an explicit consequence of the general M-018
criterion, not a stronger existential criterion.

\begin{proof}
Take $\gamma\in\Gamma_\lambda$. Remove at most $b$
variable factors, always starting with the largest index,
and call their multiset $\eta$. If $|\gamma|\leq b$,
remove all factors. Set $\alpha=\gamma-\eta$.
For every suffix beginning at $j$, this procedure removes
exactly the smaller of $b$ and that suffix's original
size before removing any earlier variable. Consequently
\[
         \sum_{i=j}^r\alpha_i
                =\left(\sum_{i=j}^r\gamma_i-b\right)_+
                \leq(\lambda_j-b)_+=\mu_j.
                                                               \tag{O-024.3}
\]
Hence $\alpha\in\Gamma_\mu$ and $\eta\in D_b$,
proving (O-024.1).

For (O-024.2), send $\alpha\in\Gamma_\mu$ to its
suffix cumulative sums
$\nu_j=\sum_{i=j}^r\alpha_i$.
They form a weakly decreasing sequence of nonnegative
integers bounded componentwise by $\mu$.
Conversely any such subpartition gives
$\alpha_j=\nu_j-\nu_{j+1}$, with $\nu_{r+1}=0$.
These two operations are inverse, proving the positive count.
\end{proof}

For an even jet order $b=2\ell+2$, put
\[
 T=\lceil\sqrt{r-1}\rceil\binom{r-1+\ell}{\ell},
 \qquad c(\mu)=\#\{\nu:\nu\subseteq\mu\}.
\]
M-018 now gives the explicit sufficient condition
\[
 n\geq T(1+c(\mu))
       \quad\Longrightarrow\quad K_n(\lambda)=a_\lambda(d[n]).
                                                               \tag{O-024.4}
\]
If $n\geq m\geq T(1+c(\mu))$, it also gives
$q^m_\lambda=K_n=a_\lambda$ with internal padding.
These are full-copy conclusions at arbitrary outer degree.

\paragraph{An exact positive recurrence.}
The count requires no cancellation. For $0\leq t\leq\mu_r$
put $F_r(t)=1$, and put it zero otherwise.
Recursively define
\[
 F_j(t)=\mathbf1_{\{0\leq t\leq\mu_j\}}
                         \sum_{s=0}^t F_{j+1}(s)
                    \quad(j=r-1,\ldots,2).
                                                               \tag{O-024.5}
\]
Here $F_j(t)$ counts allowed suffixes with cumulative
sum $t$ at position $j$; choosing the next cumulative
sum $s\leq t$ proves the recurrence.
Finally $c(\mu)=\sum_{t=0}^{\mu_2}F_2(t)$.
This is a finite positive computation, not a claim of
polynomial complexity in a binary encoding of the row lengths.

\paragraph{A sharper explicit three-row condition in BIP range.}
Every actual obstruction at $n\geq m^{25}$ satisfies
\[
 \boxed{\lambda_2\geq159375611\quad\text{or}\quad
          \lambda_3\geq29\quad\text{or}\quad\lambda_4\geq27.}
                                                               \tag{O-024.6}
\]
Indeed suppose the three inequalities fail. With $b=26$,
the shaved tail has at most two nonzero parts, bounded
by $(159375584,2)$. If a two-part partition is $(A,C)$,
with $A\geq C$, its number of subpartitions is
\[
                  \sum_{j=0}^C(A-j+1)
                 =(C+1)(A+1)-\frac{C(C+1)}2.
                                                               \tag{O-024.7}
\]
Thus here
$c(\mu)\leq478126752$ and
$1+c(\mu)\leq478126753<C_*=478126755$.
The uniform normalized estimate of M-017 makes
$T(1+c(\mu))<m^{25}\leq n$, so (O-024.4)
gives $K_n=a_\lambda$ and excludes an actual obstruction.
The exact arithmetic also gives
$3\cdot159375611-77=478126756>C_*$;
the displayed constant is the largest obtained by this
particular two-part count with the stated other caps.
It is not optimal among all order-ideal choices.

\paragraph{The canonical order ideal need not be smallest.}
For the full degree-four simplex in two variables, with
$b=2$, the canonical shaved ideal is the degree-two
simplex and has six members. The five-member order ideal
\[
                     \{1,x,x^2,y,y^2\}
\]
also covers the degree-four simplex after adding $D_2$:
a monomial of degree at most two uses the member $1$;
one of degree three or four has some exponent at least
two, so remove $x^2$ or $y^2$.
Thus (O-024.2) is an explicit sufficient count, not
the solution of the finite optimization in M-018.

\paragraph{Grothendieck dictionary.}
The exponent sets are finite order ideals in a torus
character monoid, and their free coefficient modules
are comodules as in SGA~3 I, Proposition~4.7.3,
pp.~35--36. The positive count (O-024.2) indexes
an explicit integral monomial basis. M-018 supplies
the actual universally injective coordinate pullback;
F-016's Schur/exterior support calculation uses SGA~6
VI, Theorem~3.3, pp.~392--393, and Macdonald I(4.3),
p.~63, Appendix~A(5.3)--(5.6), pp.~155--156.
The partition bijection gives the size of a coefficient
module. It does not by itself give positivity of a new
Kronecker coefficient or saturation of its support;
the multiplicity conclusion still uses the proved
dominant coefficient morphism.
